% ============================================================================
% 第一章：研究背景与意义
% ============================================================================

\section{研究背景}

\begin{frame}{研究背景}
    \begin{block}{运动想象脑机接口 (MI-BCI)}
        \begin{itemize}
            \item 通过解码运动想象诱发的 ERD 现象实现用户意图识别
            \item 为运动功能障碍患者提供重建运动功能的可能性
            \item 无需外部刺激、用户可自主控制
        \end{itemize}
    \end{block}
    
    \vspace{0.5cm}
    \begin{alertblock}{关键问题：时间窗选择}
        \begin{itemize}
            \item ERD 动态特性存在显著的\textbf{个体间差异}
            \item 时间窗长度需在\textbf{信息完整性}与\textbf{噪声抑制}之间平衡
            \item 在线 BCI 系统需要\textbf{毫秒级响应}
        \end{itemize}
    \end{alertblock}
\end{frame}

\begin{frame}{时间窗选择的重要性}
    \begin{center}
    \begin{tabular}{|c|c|c|}
        \hline
        \textbf{因素} & \textbf{影响} & \textbf{挑战} \\
        \hline
        个体差异 & ERD 模式不同 & 固定窗口无法适应 \\
        \hline
        信噪比 & 窗口长度影响 & 需要权衡 \\
        \hline
        实时性 & 响应延迟 & 在线决策要求 \\
        \hline
    \end{tabular}
    \end{center}
    
    \vspace{0.5cm}
    \begin{exampleblock}{结论}
        自适应时间窗优化是提升 MI-BCI 分类性能的关键途径
    \end{exampleblock}
\end{frame}

% ============================================================================
% 问题提出
% ============================================================================
\section{问题提出}

\begin{frame}{现有方法及其局限}
    \begin{table}[h]
        \centering
        \footnotesize
        \begin{tabular}{l|c|c|c}
            \hline
            \textbf{方法} & \textbf{计算开销} & \textbf{适应性} & \textbf{可解释性} \\
            \hline
            固定时间窗 & 低 & 差 & 好 \\
            滑动窗 & 中 & 中 & 中 \\
            贝叶斯优化 & 高 & 中 & 中 \\
            深度强化学习 & 很高 & 好 & 差 \\
            \hline
        \end{tabular}
    \end{table}
    
    \vspace{0.5cm}
    \begin{alertblock}{深度强化学习的困境}
        \begin{itemize}
            \item 百万级参数 $\rightarrow$ 难以嵌入式部署
            \item 训练不稳定 $\rightarrow$ 对超参数敏感
            \item 黑盒决策 $\rightarrow$ 难以调试分析
            \item 数据需求大 $\rightarrow$ 与 BCI 数据稀缺矛盾
        \end{itemize}
    \end{alertblock}
\end{frame}

\begin{frame}{核心研究问题}
    \begin{alertblock}{研究问题}
        在运动想象时间窗优化任务中，是否存在一种方法，能够在保持与深度强化学习相当性能的同时，显著降低模型复杂度、提升训练稳定性和可解释性？
    \end{alertblock}
    
    \vspace{0.5cm}
    \begin{exampleblock}{核心矛盾}
        \begin{center}
        深度强化学习对大量训练数据和高计算资源的依赖 \\
        \vspace{0.3cm}
        \Large\textbf{vs} \\
        \vspace{0.3cm}
        BCI 系统中数据稀缺、设备资源受限的现实
        \end{center}
    \end{exampleblock}
\end{frame}

% ============================================================================
% 研究内容
% ============================================================================
\section{研究内容}

\begin{frame}{研究思路}
    \begin{block}{核心洞察}
        时间窗优化任务的状态空间是\textbf{有限的}（136 个状态）
    \end{block}
    
    \vspace{0.5cm}
    \begin{exampleblock}{关键问题}
        在状态空间有限的特定任务中，是否真的需要复杂的深度神经网络？
    \end{exampleblock}
    
    \vspace{0.5cm}
    \begin{alertblock}{本研究方案}
        使用\textbf{表格型 Q-Learning}替代深度 Q 网络
        \begin{itemize}
            \item 维护状态 - 动作价值表（Q 表）
            \item 无需函数近似
            \item 理论收敛保证
        \end{itemize}
    \end{alertblock}
\end{frame}

\begin{frame}{研究目标}
    \begin{block}{总体目标}
        设计一种轻量化的时间窗优化方法，在保持性能的同时满足便携式 BCI 设备的资源约束
    \end{block}
    
    \vspace{0.5cm}
    \begin{itemize}
        \item \textbf{性能目标}: 与深度强化学习方法相当（准确率差异不显著）
        \item \textbf{效率目标}: 显著降低计算资源需求（训练时间、推理延迟、内存占用）
        \item \textbf{可解释性目标}: 决策过程可追踪、可分析
    \end{itemize}
\end{frame}

\begin{frame}{技术路线}
    \begin{enumerate}
        \item \textbf{问题建模}
        \begin{itemize}
            \item 将时间窗优化形式化为马尔可夫决策过程 (MDP)
            \item 定义状态空间、动作空间、奖励函数
        \end{itemize}
        
        \item \textbf{算法设计}
        \begin{itemize}
            \item 设计表格型 Q-Learning 算法
            \item 采用$\varepsilon$-贪婪策略平衡探索与利用
        \end{itemize}
        
        \item \textbf{实验验证}
        \begin{itemize}
            \item 在 BCI Competition IV 2a 数据集上评估
            \item 与多种基线方法进行对比分析
        \end{itemize}
    \end{enumerate}
\end{frame}

% ============================================================================
% 主要贡献
% ============================================================================
\section{主要贡献}

\begin{frame}{主要贡献}
    \begin{enumerate}
        \item \textbf{轻量化设计范式}
        \begin{itemize}
            \item 将深度强化学习简化为表格型强化学习
            \item 544 个参数 vs DQN 百万级参数
        \end{itemize}
        
        \item \textbf{算法选择准则}
        \begin{itemize}
            \item 状态空间 $< 10^4$ 时，表格型方法优于 DQN
            \item 收敛性保证、样本效率高、可解释性强
        \end{itemize}
        
        \item \textbf{实验验证}
        \begin{itemize}
            \item BCI Competition IV 2a 数据集
            \item 性能与 DQN 相当 ($p = 0.981$)
            \item 训练加速 3 倍、推理延迟降低 1000 倍
        \end{itemize}
    \end{enumerate}
\end{frame}

\begin{frame}{核心实验结果}
    \begin{table}[h]
        \centering
        \begin{tabular}{lcc}
            \hline
            \textbf{指标} & \textbf{Table Q} & \textbf{DQN} \\
            \hline
            准确率 & 62.68\% & 62.70\% \\
            统计检验 & \multicolumn{2}{c}{$p = 0.981$} \\
            训练时间 & \textbf{快 3 倍} & 慢 \\
            推理延迟 & \textbf{低 1000 倍} & 高 \\
            内存占用 & \textbf{低 125 倍} & 高 \\
            收敛性 & \checkmark 理论保证 & $\times$ 无保证 \\
            可解释性 & \checkmark 完全可解释 & $\times$ 黑盒 \\
            \hline
        \end{tabular}
    \end{table}
    
    \vspace{0.5cm}
    \begin{alertblock}{核心结论}
        表格型 Q-Learning 在保持与 DQN 相当性能的同时，显著降低了计算资源需求
    \end{alertblock}
\end{frame}

\begin{frame}{本章小结}
    \begin{block}{研究背景}
        时间窗选择是 MI-BCI 性能的关键影响因素
    \end{block}
    
    \vspace{0.3cm}
    \begin{block}{核心问题}
        如何在保持性能的前提下实现轻量化、可解释的时间窗优化
    \end{block}
    
    \vspace{0.3cm}
    \begin{block}{研究方案}
        使用表格型 Q-Learning 替代深度 Q 网络
    \end{block}
    
    \vspace{0.3cm}
    \begin{block}{主要贡献}
        轻量化设计范式、算法选择准则、实验验证
    \end{block}
\end{frame}
